\begin{problem}[第35届全国中学生物理竞赛决赛][平行板电容器与导体板运动]{97}
    平行板电容器极板 1 和 2 的面积均为 $S$，水平固定放置，它们之间的距离为 $d$，接入如图所示的电路中，电源的电动势记为 $U$。不带电的导体薄平板 3 的质量为 $m$、尺寸与电容器极板相同。平板 3 平放在极板 2 的正上方，且与极板 2 有良好的电接触。整个系统置于真空室内，真空的介电常量为 $\varepsilon_{0}$ 。闭合电键 K 后，平板 3 与极板 1 和 2 相继碰撞，上下往复运动。假设导体板之间的电场均可视为匀强电场；导线电阻和电源内阻足够小，充放电时间可忽略不计；平板 3 与极板 1 或 2 碰撞后立即在极短时间内达到静电平衡；所有碰撞都是完全非弹性的。重力加速度大小为 $g$ 。
    \begin{subproblem}
        电源电动势 $U$ 至少为多大?
    \end{subproblem}
    \begin{subproblem}
        求平板3运动的周期（用 $U$ 和题给条件表示）。

        已知积分公式
        $\int \frac{\mathrm{d}x}{\sqrt{ax^2 + bx}} = \frac{1}{\sqrt{a}}\ln \left(2ax + b + 2\sqrt{a}\sqrt{ax^2 + bx}\right) + C$ ，其中 $a > 0$ ， $C$ 为积分常数。
    \end{subproblem}
\end{problem}